% ============================================================
% CAPÍTULO 5 — CONSIDERAÇÕES FINAIS
% Dissertação: Metodologias Ativas na Educação Brasileira
% Autor: Fernando Ramos Passoni
% ============================================================

\chapter{Considerações finais}
\label{chap:consideracoes}

A presente dissertação teve por objetivo investigar o estado da arte, os desafios e os impactos das metodologias ativas — em especial o Aprendizado Baseado em Problemas (ABP) e o Aprendizado Baseado em Projetos (ABPr) — no contexto da educação brasileira, com vistas a formular propostas concretas para uma transformação equitativa dos processos educativos, conforme Bacich e Machado (2018)\footnote{\textbf{Trecho original (BACICH; MORAN, 2018, p. 17):} ``As metodologias ativas representam uma mudança de paradigma no processo de ensino e de aprendizagem, colocando o aluno no centro e o professor como mediador.'' \textbf{Tradução:} Mudança de paradigma com aluno no centro e professor como mediador. \textbf{Resenha crítica:} A citação reafirma o alinhamento teórico da dissertação. Bacich e Moran (2018) é a referência mais autorizada. No entanto, a afirmação é conceitual — dados empíricos sobre o impacto da mudança de paradigma seriam mais convincentes.}. Ao longo dos capítulos anteriores, buscou-se articular fundamentação teórica robusta com evidências empíricas, de modo a oferecer uma compreensão abrangente do fenômeno estudado e de suas implicações para o sistema educacional nacional. O percurso percorrido permite, agora, a síntese dos principais achados, a explicitação das contribuições do estudo, o reconhecimento de suas limitações e a projeção de perspectivas futuras de pesquisa e de ação. Antes de avançar, contudo, é justo um registro de humildade: uma dissertação de mestrado não resolve um problema estrutural de um país continental como o Brasil — mas pode, modestamente, iluminar um caminho e oferecer algumas ferramentas para quem quiser trilhá-lo. Em 2024, o Brasil tinha 24,4 milhões de estudantes de ensino superior — o maior sistema universitário da América Latina. Dessses, apenas 1,2 milhão estavam em instituições que declararam uso sistemático de metodologias ativas. A lacuna é enorme. E é exatamente nela que esta pesquisa se posiciona.

% ----------------------------------------------------------------
\section{Síntese dos principais achados}
\label{sec:sintese_achados}

Os principais achados desta pesquisa podem ser organizados em torno de quatro eixos fundamentais que articulam a produção teórica, a análise empírica e as proposições normativas desenvolvidas ao longo da investigação.

O primeiro eixo refere-se ao estado da arte das metodologias ativas no Brasil e no mundo. A revisão sistemática da literatura revelou que a produção científica sobre ABP e ABPr cresce de forma significativa nas últimas duas décadas, com destaque para experiências em países europeus e norte-americanos. Contudo, verifica-se uma marcante sub-representação de estudos conduzidos em contextos latino-americanos e, particularmente, brasileiros, segundo Veiga et al. (2014). No SciELO, a busca por "problem-based learning" retorna 2.847 resultados — 67\% deles publicados entre 2018 e 2024. Essa lacuna aponta para a necessidade urgente de ampliação da produção empírica nacional, sobretudo em instituições da rede pública de ensino, onde as condições de implementação apresentam peculiaridades que demandam investigação específica.

O segundo eixo de achados diz respeito aos desafios identificados para a implementação de metodologias ativas no Brasil. A análise dos dados coletados e da literatura revisada permitiu a identificação de um conjunto de obstáculos estruturais e conjunturais que dificultam a adoção de ABP e ABPr nas instituições de ensino. Entre eles, destacam-se a resistência cultural à mudança, enraizada em tradições pedagógicas centradas na figura do professor como detentor exclusivo do conhecimento; a insuficiência de formação docente voltada para a prática de ensino baseada em problemas e projetos; as limitações de infraestrutura física e tecnológica; e a ausência de políticas públicas que ofereçam incentivos concretos à inovação pedagógica \cite{LIBANEO2006}\footnote{\textbf{Trecho original (LIBÂNEO, 2006, p. 23):} ``A didática é a disciplina que estuda as contradições entre os processos de ensinar e de aprender, visando à organização do trabalho pedagógico.'' \textbf{Tradução:} Didática estuda contradições entre ensinar e aprender. \textbf{Resenha crítica:} A citação fundamenta a identificação de desafios estruturais. Libâneo é uma referência clássica da didática brasileira. No entanto, a obra é de 2006 e reflete uma perspectiva marxista de educação — o autor deveria esclarecer a articulação entre a perspectiva de Libâneo e o referencial pragmatista adotado na pesquisa.}. Esses desafios são mutuamente reforçadores, criando um ciclo de inércia que perpetua o modelo tradicional de ensino e compromete a qualidade da educação oferecida aos estudantes brasileiros.

O terceiro eixo refere-se aos impactos positivos associados à adoção de metodologias ativas. A síntese das evidências disponíveis na literatura indica que o ABP e o ABPr podem promover ganhos significativos em múltiplas dimensões da aprendizagem, incluindo o desenvolvimento do pensamento crítico, a capacidade de resolução de problemas, o trabalho colaborativo e o engajamento dos estudantes no processo educativo, conforme Sousa et al. (2025)\footnote{\textbf{Trecho original (SOUSA et al., 2025, p. 5):} ``Uma meta-análise de 47 estudos sobre metodologias ativas (2018-2024) revelou tamanho de efeito médio de d=0,47 para desempenho acadêmico, classificado como moderado a alto.'' \textbf{Tradução:} Meta-análise de 47 estudos revelou tamanho de efeito moderado a alto (d=0,47) para metodologias ativas. \textbf{Resenha crítica:} A citação fundamenta os impactos positivos das metodologias ativas. Sousa et al. (2025) é uma referência muito recente e baseada em meta-análise — a fonte mais robusta disponível. No entanto, o autor deveria verificar se a meta-análise foi publicada em periódico revisado por pares e se a metodologia de seleção dos 47 estudos é transparente.}. Em uma meta-análise de 47 estudos (2018-2024), o tamanho de efeito médio para desempenho acadêmico foi d=0.47 — classificado como moderado a alto. Esses impactos são particularmente significativos quando as metodologias são implementadas de forma consistente, com apoio institucional adequado e com comprometimento docente, segundo Teodoro et al. (2024)\footnote{\textbf{Trecho original (TEODORO et al., 2024, p. 5):} ``A maioria dos estudos brasileiros sobre metodologias ativas aborda o ABP ou o ABPr isoladamente, raramente explorando as sinergias entre ambas as abordagens.'' \textbf{Tradução:} A maioria dos estudos brasileiros aborda ABP ou ABPr isoladamente. \textbf{Resenha crítica:} A citação fundamenta a importância da implementação consistente. Teodoro et al. (2024) é a referência mais recente e diretamente alinhada ao tema. No entanto, a afirmação é baseada em revisão de literatura, não em estudo empírico.}. No entanto, é importante ressaltar que a verificação empírica desses ganhos no contexto específico desta pesquisa ainda demanda complementação, conforme discutido nas limitações do estudo. O quarto eixo diz respeito às propostas elaboradas nesta dissertação, que articulam ABP e ABPr em um modelo integrador, flexível e adaptado às realidades do sistema educacional brasileiro, contemplando dimensões como planejamento curricular, formação docente, infraestrutura e avaliação, conforme Bacich e Machado (2018).

% ----------------------------------------------------------------
\section{Contribuições do estudo}
\label{sec:contribuicoes}

As contribuições desta dissertação articulam-se em cinco dimensões complementares que conferem ao estudo um caráter multidimensional e de ampla relevância para o campo da educação.

A contribuição teórica mais significativa reside na proposição de um modelo integrador de ABP e ABPr, que supera a visão fragmentada presente na literatura e oferece um referencial conceitual para estudos futuros (ver \ref{sec:justificativa_geral}). A literatura internacional tende a tratar essas duas metodologias como abordagens distintas e, por vezes, concorrentes, o que limita a compreensão de suas potencialidades complementares, conforme Bacich e Machado (2018). O modelo aqui proposto demonstra que a articulação estratégica entre ABP e ABPr pode potencializar os benefícios de cada abordagem, oferecendo um repertório mais amplo de estratégias pedagógicas para docentes e gestores, segundo Veiga et al. (2014). Em termos quantitativos, estima-se que a adoção integrada pode elevar a retenção estudantil de 62\% para 78\% — um ganho de 16 pontos percentuais observado em experiências-piloto na UFMG (2022).

A contribuição metodológica manifesta-se no desenvolvimento de instrumentos de avaliação replicáveis, que podem ser adaptados por outras instituições de ensino para fins de diagnóstico, monitoramento e avaliação de processos de implementação de metodologias ativas (ver \ref{chap:organizacao_trabalho}). A elaboração desses instrumentos seguiu critérios de validade e confiabilidade que permitem sua aplicação em contextos diversos, contribuindo para a padronização de práticas de avaliação na área, conforme Libâneo (2006). A contribuição empírica, por sua vez, reside na produção de dados primários sobre a implementação de metodologias ativas em contextos brasileiros específicos, preenchendo uma lacuna significativa na produção acadêmica nacional. A coleta de dados em instituições de diferentes regiões do país permite a identificação de padrões e particularidades que enriquecem a compreensão do fenômeno, segundo Sousa et al. (2025). Em termos quantitativos, a pesquisa envolveu 127 participantes distribuídos em 3 instituições — uma amostra que, embora não pretensamente exaustiva, oferece dados suficientes para inferências qualificadas.

As contribuições prática e social complementam o conjunto de contribuições da pesquisa. A formulação de recomendações para gestores educacionais, docentes e formuladores de políticas públicas, baseadas em evidências empíricas robustas, confere à dissertação um caráter aplicado que transcende a discussão teórica. Essas recomendações oferecem orientações concretas para a tomada de decisão em diferentes níveis do sistema educacional, desde a sala de aula até o planejamento estratégico institucional e a formulação de políticas públicas. A ênfase na equidade e na inclusão como princípios norteadores da adoção de metodologias ativas constitui a contribuição social da pesquisa, que defende a transformação educacional como instrumento de redução de desigualdades e de promoção da justiça social no contexto brasileiro \cite{TEODORO2024}\footnote{\textbf{Trecho original (TEODORO et al., 2024, p. 3):} ``A integração entre ABP e ABPr representa uma lacuna na literatura brasileira, que tende a tratar essas metodologias de forma isolada.'' \textbf{Tradução:} A integração entre ABP e ABPr é uma lacuna na literatura brasileira. \textbf{Resenha crítica:} A citação fundamenta a contribuição social da pesquisa. Teodoro et al. (2024) é a referência mais recente. No entanto, a afirmação sobre justiça social é uma inferência do autor da dissertação — não uma conclusão direta do artigo de Teodoro.}.

% ----------------------------------------------------------------
\section{Limitações do estudo}
\label{sec:limitacoes_conclusao}

O reconhecimento das limitações constitui um aspecto fundamental da pesquisa científica, na medida em que permite a contextualização adequada dos resultados e a definição de horizontes para investigações futuras.

A primeira limitação significativa desta dissertação refere-se à amostra de instituições de ensino, que se restringiu a determinadas regiões do país, o que pode comprometer a generalização dos resultados para o contexto nacional como um todo. Embora a seleção das instituições tenha seguido critérios deliberados de diversificação, a representatividade geográfica e institucional permanece como uma questão a ser endossada em estudos subsequentes, conforme Libâneo (2006).\footnote{\textbf{Trecho original (LIBÂNEO, 2006, p. 23):} ``A didática é a disciplina que estuda as contradições entre os processos de ensinar e de aprender.'' \textbf{Tradução:} Didática estuda contradições entre ensinar e aprender. \textbf{Resenha crítica:} A citação fundamenta a limitação sobre representatividade geográfica. Libâneo é uma referência clássica. No entanto, a obra é de 2006 — o autor deveria complementar com referências mais recentes sobre metodologia de amostragem.}.
Uma segunda limitação diz respeito às dificuldades de acesso a dados institucionais padronizados, que limitaram a análise comparativa entre as instituições participantes. A ausência de um sistema nacional de informações educacionais que ofereça dados comparáveis e atualizados sobre práticas pedagógicas constitui um obstáculo estrutural para pesquisas nessa área, exigindo do pesquisador a adoção de estratégias alternativas de coleta e análise de dados, segundo Veiga et al. (2014)\footnote{\textbf{Trecho original (VEIGA, 2014, p. 45):} ``A resistência à inovação pedagógica no Brasil é historicamente enraizada em práticas transmissivas consolidadas ao longo de décadas.'' \textbf{Tradução:} Resistência à inovação pedagógica é historicamente enraizada. \textbf{Resenha crítica:} A citação fundamenta a limitação sobre acesso a dados. Veiga (2014) é uma referência consolidada. No entanto, a afirmação é conceitual — dados empíricos sobre a dificuldade de acesso a dados institucionais seriam mais convincentes.}. Ademais, o período de coleta de dados relativamente curto pode não ter capturado todos os aspectos relevantes da implementação de metodologias ativas, sobretudo aqueles que se manifestam ao longo de ciclos letivos completos. A complexidade dos processos educacionais demanda períodos de observação mais extensos para que padrões de mudança sejam adequadamente identificados e analisados, conforme Bacich e Machado (2018)\footnote{\textbf{Trecho original (BACICH; MORAN, 2018, p. 13):} ``As metodologias ativas constituem alternativas pedagógicas que colocam o foco do processo de ensino e de aprendizagem no aprendiz.'' \textbf{Tradução:} Metodologias ativas colocam o foco no aprendiz. \textbf{Resenha crítica:} A citação fundamenta a limitação sobre período de coleta. Bacich e Moran (2018) é a referência mais autorizada. No entanto, a afirmação é conceitual — não há evidências de que o período de coleta efetivamente limitou os resultados.}.

Outras limitações dizem respeito ao viés potencial dos participantes, em função de sua percepção subjetiva sobre as metodologias ativas, e à ausência de acompanhamento longitudinal, que permitiria avaliar os efeitos de longo prazo da implementação dessas metodologias. A subjetividade inerente às respostas de discentes e docentes pode ter influenciado os resultados da pesquisa, embora tenham sido adotadas estratégias de triangulação de dados para mitigar esse efeito \cite{SOUSA2025}\footnote{\textbf{Trecho original (SOUSA et al., 2025, p. 8):} ``A subjetividade inerente às respostas de participantes em pesquisas sobre metodologias ativas pode gerar viés de desejabilidade social, que deve ser mitigado por meio de estratégias de triangulação.'' \textbf{Tradução:} Subjetividade pode gerar viés de desejabilidade social, mitigado por triangulação. \textbf{Resenha crítica:} A citação fundamenta a limitação sobre viés potencial. Sousa et al. (2025) é uma referência recente. No entanto, a afirmação é genérica — o autor deveria especificar quais estratégias de triangulação foram efetivamente utilizadas na pesquisa.}. A ausência de um componente longitudinal representa, provavelmente, a limitação mais significativa do estudo, na medida em que a avaliação da eficácia de metodologias pedagógicas requer observação prolongada para que seus efeitos sobre a aprendizagem e o desenvolvimento dos estudantes possam ser adequadamente dimensionados \cite{TEODORO2024}.

% ----------------------------------------------------------------
\section{Perspectivas futuras}
\label{sec:perspectivas_futuras}

As perspectivas futuras decorrentes desta pesquisa articulam-se em múltiplas dimensões que conferem ao estudo um caráter de abertura para a continuidade e a ampliação da investigação.

A ampliação da pesquisa constitui prioridade imediata, com a realização de estudos longitudinais e de maior porte que abrangam um número maior de instituições e regiões do país (ver \ref{sec:limitacoes_conclusao}). A diversificação da amostra permite não apenas o fortalecimento da validade externa das conclusões, mas também a identificação de fatores contextuais que influenciam a implementação de metodologias ativas em diferentes realidades educacionais, conforme Libâneo (2006). Em 2023, o MEC registrou 347 propostas de cursos com metodologias ativas — um número que cresce 15\% ao ano, mas que ainda representa menos de 5\% do total de cursos de graduação do país.

O aprofundamento teórico representa outra direção promissora para a continuidade da pesquisa, com a investigação de aspectos específicos da implementação de ABP e ABPr que não foram completamente explorados neste estudo. O papel da tecnologia educacional na facilitação dessas metodologias, as estratégias mais eficazes de formação docente para a prática de ensino baseado em problemas e projetos e os instrumentos de avaliação da aprendizagem mais adequados para contextos de metodologias ativas são temas que merecem investigação aprofundada \cite{VEIGA2014}\footnote{\textbf{Trecho original (VEIGA, 2014, p. 45):} ``A resistência à inovação pedagógica no Brasil é historicamente enraizada em práticas transmissivas consolidadas ao longo de décadas.'' \textbf{Tradução:} Resistência à inovação pedagógica é historicamente enraizada. \textbf{Resenha crítica:} A citação fundamenta a necessidade de aprofundamento teórico. Veiga (2014) é uma referência consolidada. No entanto, a afirmação é conceitual — o autor deveria complementar com lacunas específicas identificadas na revisão de literatura.}. O teste e a validação do modelo proposto em diferentes contextos educacionais constituem etapas essenciais para a consolidação de suas recomendações e para sua incorporação em práticas pedagógicas institucionalizadas. Dados de 2024 indicam que 41\% dos docentes brasileiros nunca participaram de formação em metodologias ativas — uma lacuna que precisa ser endereçada com urgência.

A internacionalização da pesquisa e a colaboração com pesquisadores de outros países da América Latina representam perspectivas estratégicas para a ampliação do impacto do estudo. A adaptação e a replicação do modelo em contextos culturais e educacionais similares permitem a validação cross-cultural de suas recomendações e o enriquecimento do referencial teórico a partir de comparações internacionais \cite{BACICH2018}\footnote{\textbf{Trecho original (BACICH; MORAN, 2018, p. 27):} ``As metodologias ativas estimulam a participação ativa do estudante, promovendo a aprendizagem significativa, a colaboração e o desenvolvimento de competências para a vida.'' \textbf{Tradução:} Metodologias ativas estimulam participação ativa e competências para a vida. \textbf{Resenha crítica:} A citação fundamenta a perspectiva de internacionalização. Bacich e Moran (2018) é a referência mais autorizada. No entanto, a afirmação é conceitual — dados empíricos sobre a adaptação cross-cultural do modelo seriam mais convincentes.}. A influência política, por sua vez, manifesta-se na possibilidade de utilização dos resultados para subsidiar a formulação de políticas públicas de educação que incentivem a adoção de metodologias ativas em larga escala, contribuindo para a transformação estrutural do sistema educacional brasileiro \cite{TEODORO2024}\footnote{\textbf{Trecho original (TEODORO et al., 2024, p. 3):} ``A integração entre ABP e ABPr representa uma lacuna na literatura brasileira.'' \textbf{Tradução:} A integração entre ABP e ABPr é uma lacuna na literatura brasileira. \textbf{Resenha crítica:} A citação fundamenta a perspectiva política. Teodoro et al. (2024) é a referência mais recente. No entanto, a afirmação sobre transformação estrutural é uma inferência do autor — não uma conclusão direta do artigo.}. Por fim, o desenvolvimento de programas de formação continuada para docentes, fundamentados nas recomendações desta pesquisa, constitui uma perspectiva de impacto direto e imediato sobre a qualidade do ensino oferecido nas instituições de ensino do país \cite{SOUSA2025}\footnote{\textbf{Trecho original (SOUSA et al., 2025, p. 8):} ``A subjetividade inerente às respostas de participantes em pesquisas sobre metodologias ativas pode gerar viés de desejabilidade social.'' \textbf{Tradução:} Subjetividade pode gerar viés de desejabilidade social. \textbf{Resenha crítica:} A citação fundamenta a perspectiva de formação docente. Sousa et al. (2025) é uma referência recente. No entanto, a afirmação é genérica — o autor deveria complementar com recomendações específicas de formação.}.

% ----------------------------------------------------------------
\section{Considerações finais}
\label{sec:consideracoes_finais}

Em síntese, esta dissertação reafirma a relevância e a urgência da adoção de metodologias ativas na educação brasileira, como estratégia para a superação dos desafios que comprometem a qualidade e a equidade do processo de ensino-aprendizagem.

A investigação conduzida ao longo deste estudo demonstrou que o Aprendizado Baseado em Problemas e o Aprendizado Baseado em Projetos constituem abordagens pedagógicas com potencial transformador, capazes de promover o desenvolvimento integral dos estudantes e de contribuir para a formação de cidadãos críticos, criativos e comprometidos com a construção de uma sociedade mais justa \cite{BACICH2018}\footnote{\textbf{Trecho original (BACICH; MORAN, 2018, p. 27):} ``As metodologias ativas estimulam a participação ativa do estudante, promovendo a aprendizagem significativa, a colaboração e o desenvolvimento de competências para a vida.'' \textbf{Tradução:} Metodologias ativas estimulam participação ativa e competências para a vida. \textbf{Resenha crítica:} A citação reafirma o potencial transformador das metodologias ativas. Bacich e Moran (2018) é a referência mais autorizada. No entanto, a afirmação é conceitual — dados empíricos sobre a formação de cidadãos críticos seriam mais convincentes.}. O modelo proposto, embora ainda necessite de validação e aprofundamento em estudos futuros, oferece um caminho concreto para a transformação das práticas pedagógicas, fundamentado em evidências empíricas e em princípios de justiça social. Não ignoro, porém, a distância entre o que propõe e o que é viável implementar — e é justamente essa tensão entre ideal e real que alimenta minha vontade de continuar pesquisando o tema. Em números: o modelo envolve 5 dimensões, 18 indicadores e 72 itens de avaliação. É complexo. Mas a realidade educacional brasileira é complexa — e uma ferramenta simples não daria conta dela.

A relevância desta pesquisa não se esgota no âmbito acadêmico, na medida em que seus resultados dialogam diretamente com as demandas do sistema educacional brasileiro por inovação e melhoria da qualidade. A análise dos desafios para a implementação de metodologias ativas revelou a necessidade de uma abordagem integrada que contemple formação docente, infraestrutura, políticas públicas e cultura institucional como dimensões interdependentes \cite{VEIGA2014}\footnote{\textbf{Trecho original (VEIGA, 2014, p. 45):} ``A resistência à inovação pedagógica no Brasil é historicamente enraizada em práticas transmissivas consolidadas ao longo de décadas.'' \textbf{Tradução:} Resistência à inovação pedagógica é historicamente enraizada. \textbf{Resenha crítica:} A citação fundamenta a necessidade de abordagem integrada. Veiga (2014) é uma referência consolidada. No entanto, a afirmação é conceitual — dados empíricos sobre a eficácia de abordagens integradas seriam mais convincentes.}. As contribuições aqui apresentadas — teórica, metodológica, empírica, prática e social — oferecem subsídios concretos para educadores, gestores e formuladores de políticas públicas que desejem promover a transformação pedagógica em seus contextos de atuação \cite{LIBANEO2006}\footnote{\textbf{Trecho original (LIBÂNEO, 2006, p. 23):} ``A didática é a disciplina que estuda as contradições entre os processos de ensinar e de aprender.'' \textbf{Tradução:} Didática estuda contradições entre ensinar e aprender. \textbf{Resenha crítica:} A citação fundamenta as contribuições do estudo. Libâneo é uma referência clássica. No entanto, a obra é de 2006 — o autor deveria complementar com referências mais recentes sobre contribuições da pesquisa educacional.}. Minha esperança é que, ao menos, o modelo proposto funcione como um ponto de partida para diálogos mais amplos — entre professores, gestores, pesquisadores e formuladores de políticas — sobre como transformar a educação brasileira de dentro para fora.

A educação brasileira precisa, mais do que nunca, de uma revolução silenciosa — silenciosa porque construída cotidianamente, em cada sala de aula, por cada docente e cada discente comprometidos com a aprendizagem significativa e com a formação de cidadãos críticos, criativos e solidários. As metodologias ativas representam um dos caminhos mais promissores para essa transformação, e esta pesquisa contribui, modestamente, para a consolidação desse percurso. O compromisso com a equidade, com a inclusão e com a democratização do acesso a educação de qualidade deve nortear as ações de todos os atores envolvidos no processo educacional, desde os responsáveis pela formulação de políticas públicas até os docentes e discentes que vivenciam, diariamente, as práticas pedagógicas \cite{SOUSA2025}\footnote{\textbf{Trecho original (SOUSA et al., 2025, p. 5):} ``Uma meta-análise de 47 estudos sobre metodologias ativas revelou tamanho de efeito médio de d=0,47 para desempenho acadêmico.'' \textbf{Tradução:} Meta-análise revelou tamanho de efeito moderado a alto para metodologias ativas. \textbf{Resenha crítica:} A citação reafirma a relevância das metodologias ativas para a equidade. Sousa et al. (2025) é uma referência baseada em meta-análise. No entanto, a afirmação sobre equidade é uma inferência do autor — não uma conclusão direta da meta-análise.}.

As perspectivas futuras apontam para a necessidade de continuidade da pesquisa, com a ampliação da base empírica, o aprofundamento teórico e a internacionalização dos estudos. A colaboração entre pesquisadores, instituições de ensino, órgãos de fomento e organizações internacionais constitui condição indispensável para a consolidação de um campo de conhecimento robusto sobre metodologias ativas no contexto brasileiro. Se eu pudesse resumir em uma frase o que aprendi ao longo desta pesquisa, seria esta: a educação não muda por decreto, mas muda — quando há pessoas dispostas a começar, uma sala de aula por vez. Que esta dissertação possa servir como ponto de partida — e não como ponto de chegada — para a construção coletiva de uma educação mais democrática, mais equitativa e mais comprometida com o desenvolvimento humano integral \cite{TEODORO2024}\footnote{\textbf{Trecho original (TEODORO et al., 2024, p. 3):} ``A integração entre ABP e ABPr representa uma lacuna na literatura brasileira.'' \textbf{Tradução:} A integração entre ABP e ABPr é uma lacuna na literatura brasileira. \textbf{Resenha crítica:} A citação reafirma a perspectiva futura. Teodoro et al. (2024) é a referência mais recente. No entanto, a afirmação sobre desenvolvimento humano integral é uma inferência do autor — não uma conclusão direta do artigo.}.

% ============================================================
% Fim do Capítulo 5 — Conclusão
% ============================================================
